
\documentclass{article}
\usepackage{colortbl}
\usepackage{makecell}
\usepackage{multirow}
\usepackage{supertabular}

\begin{document}

\newcounter{utterance}

\centering \large Interaction Transcript for game `dond', experiment `coop', episode 1 with mistral{-}large{-}3{-}azure.
\vspace{24pt}

{ \footnotesize  \setcounter{utterance}{1}
\setlength{\tabcolsep}{0pt}
\begin{supertabular}{c@{$\;$}|p{.15\linewidth}@{}p{.15\linewidth}p{.15\linewidth}p{.15\linewidth}p{.15\linewidth}p{.15\linewidth}}
    \# & $\;$A & \multicolumn{4}{c}{Game Master} & $\;\:$B\\
    \hline

    \theutterance \stepcounter{utterance}  
    & & \multicolumn{4}{p{0.6\linewidth}}{
        \cellcolor[rgb]{0.9,0.9,0.9}{
            \makecell[{{p{\linewidth}}}]{
                \texttt{\tiny{[P1$\langle$GM]}}
                \texttt{Igrate kooperativnu pregovaračku igru u kojoj treba da se dogovorite kako da podelite skup predmeta između vas i drugog igrača.} \\
\\ 
\texttt{Pravila:} \\
\texttt{(a) Vama i drugom igraču se dodeljuje skup predmeta. Svako od vas takođe dobija tajnu funkciju vrednosti, koja predstavlja koliko vrednujete svaku vrstu predmeta.} \\
\texttt{(b) Razmenjujete poruke sa drugim igračem kako biste se dogovorili ko dobija koje predmete. Svako od vas može poslati najviše 5 poruka, ili možete ranije završiti tako što ćete u bilo kom trenutku dati tajnu ponudu.} \\
\texttt{(c) Od svakog od vas se traži da podnese tajnu ponudu koja navodi predmete koje želite, formatiranu u uglastim zagradama na sledeći način: "{[}Predlog: <broj> <naziv predmeta>, <broj> <naziv predmeta>, <...>{]}"} \\
\texttt{(d) Ako su vaše ponude komplementarne, tj. ima dovoljno predmeta da se ispune obe ponude, svakom igraču se dodeljuje rezultat zasnovan na zbiru vrednosti predmeta koje je dobio. U suprotnom, oboje dobijate nula poena.} \\
\texttt{(e) Vaš cilj je da maksimizujete zbir vašeg rezultata i rezultata drugog igrača.} \\
\\ 
\texttt{Hajde da počnemo.} \\
\\ 
\texttt{Skup dostupnih predmeta je:} \\
\\ 
\texttt{2 zvona, 1 lutka, 2 konci} \\
\\ 
\texttt{Vaša tajna funkcija vrednosti je:} \\
\\ 
\texttt{zvono: 1, lutka: 2, konac: 3} \\
\\ 
\texttt{VAŽNO: Vaše poruke, osim ako je u pitanju tajna ponuda, direktno se prosleđuju drugom igraču, zato nemojte uključivati bilo kakav odgovor na pravila niti tekst koji najavljuje vašu poruku. Da biste dali tajnu ponudu, koristite naznačeni format. Ne koristite uglaste zagrade kada komunicirate sa drugim igračem ili će to biti protumačeno kao vaša tajna ponuda.} \\
            }
        }
    }
    & \\ \\

    \theutterance \stepcounter{utterance}  
    & \multicolumn{4}{p{0.6\linewidth}}{
        \cellcolor[rgb]{0.9,0.9,0.9}{
            \makecell[{{p{\linewidth}}}]{
                \texttt{\tiny{[P1$\rangle$GM]}}
                \texttt{Pozdrav! Kako bismo podelili predmete? Ja bih voleo da razgovaramo o tome šta je svakome od nas važno.} \\
            }
        }
    }
    & & \\ \\

    \theutterance \stepcounter{utterance}  
    & & & \multicolumn{2}{p{0.3\linewidth}}{
        \cellcolor[rgb]{0.9,0.9,0.9}{
            \makecell[{{p{\linewidth}}}]{
                \texttt{\tiny{[GM$|$GM]}}
                \texttt{continue} \\
            }
        }
    }
    & & \\ \\

    \theutterance \stepcounter{utterance}  
    & & & \multicolumn{2}{p{0.3\linewidth}}{
        \cellcolor[rgb]{0.9,0.9,0.9}{
            \makecell[{{p{\linewidth}}}]{
                \texttt{\tiny{[GM$|$GM]}}
                \texttt{Pozdrav! Kako bismo podelili predmete? Ja bih voleo da razgovaramo o tome šta je svakome od nas važno.} \\
            }
        }
    }
    & & \\ \\

    \theutterance \stepcounter{utterance}  
    & & \multicolumn{4}{p{0.6\linewidth}}{
        \cellcolor[rgb]{0.9,0.9,0.9}{
            \makecell[{{p{\linewidth}}}]{
                \texttt{\tiny{[GM$\rangle$P2]}}
                \texttt{Igrate kooperativnu pregovaračku igru u kojoj treba da se dogovorite kako da podelite skup predmeta između vas i drugog igrača.} \\
\\ 
\texttt{Pravila:} \\
\texttt{(a) Vama i drugom igraču se dodeljuje skup predmeta. Svako od vas takođe dobija tajnu funkciju vrednosti, koja predstavlja koliko vrednujete svaku vrstu predmeta.} \\
\texttt{(b) Razmenjujete poruke sa drugim igračem kako biste se dogovorili ko dobija koje predmete. Svako od vas može poslati najviše 5 poruka, ili možete ranije završiti tako što ćete u bilo kom trenutku dati tajnu ponudu.} \\
\texttt{(c) Od svakog od vas se traži da podnese tajnu ponudu koja navodi predmete koje želite, formatiranu u uglastim zagradama na sledeći način: "{[}Predlog: <broj> <naziv predmeta>, <broj> <naziv predmeta>, <...>{]}"} \\
\texttt{(d) Ako su vaše ponude komplementarne, tj. ima dovoljno predmeta da se ispune obe ponude, svakom igraču se dodeljuje rezultat zasnovan na zbiru vrednosti predmeta koje je dobio. U suprotnom, oboje dobijate nula poena.} \\
\texttt{(e) Vaš cilj je da maksimizujete zbir vašeg rezultata i rezultata drugog igrača.} \\
\\ 
\texttt{Hajde da počnemo.} \\
\\ 
\texttt{Skup dostupnih predmeta je:} \\
\\ 
\texttt{2 zvona, 1 lutka, 2 konci} \\
\\ 
\texttt{Vaša tajna funkcija vrednosti je:} \\
\\ 
\texttt{zvono: 0, lutka: 2, konac: 4} \\
\\ 
\texttt{VAŽNO: Vaše poruke, osim ako je u pitanju tajna ponuda, direktno se prosleđuju drugom igraču, zato nemojte uključivati bilo kakav odgovor na pravila niti tekst koji najavljuje vašu poruku. Da biste dali tajnu ponudu, koristite naznačeni format. Ne koristite uglaste zagrade kada komunicirate sa drugim igračem ili će to biti protumačeno kao vaša tajna ponuda.} \\
\\ 
\texttt{Pozdrav! Kako bismo podelili predmete? Ja bih voleo da razgovaramo o tome šta je svakome od nas važno.} \\
            }
        }
    }
    & \\ \\

    \theutterance \stepcounter{utterance}  
    & & & \multicolumn{4}{p{0.6\linewidth}}{
        \cellcolor[rgb]{0.9,0.9,0.9}{
            \makecell[{{p{\linewidth}}}]{
                \texttt{\tiny{[GM$\langle$P2]}}
                \texttt{Zdravo! Razumem da je konac za tebe najvredniji, a lutka manje. Da li bi tebi odgovaralo da uzmeš oba konca, a ja uzmem lutku i jedno zvono? Tako bismo oboje dobili nešto vrednije za nas. Šta misliš?} \\
            }
        }
    }
     \\ \\

    \theutterance \stepcounter{utterance}  
    & & & \multicolumn{2}{p{0.3\linewidth}}{
        \cellcolor[rgb]{0.9,0.9,0.9}{
            \makecell[{{p{\linewidth}}}]{
                \texttt{\tiny{[GM$|$GM]}}
                \texttt{continue} \\
            }
        }
    }
    & & \\ \\

    \theutterance \stepcounter{utterance}  
    & & & \multicolumn{2}{p{0.3\linewidth}}{
        \cellcolor[rgb]{0.9,0.9,0.9}{
            \makecell[{{p{\linewidth}}}]{
                \texttt{\tiny{[GM$|$GM]}}
                \texttt{Zdravo! Razumem da je konac za tebe najvredniji, a lutka manje. Da li bi tebi odgovaralo da uzmeš oba konca, a ja uzmem lutku i jedno zvono? Tako bismo oboje dobili nešto vrednije za nas. Šta misliš?} \\
            }
        }
    }
    & & \\ \\

    \theutterance \stepcounter{utterance}  
    & & \multicolumn{4}{p{0.6\linewidth}}{
        \cellcolor[rgb]{0.9,0.9,0.9}{
            \makecell[{{p{\linewidth}}}]{
                \texttt{\tiny{[P1$\langle$GM]}}
                \texttt{Zdravo! Razumem da je konac za tebe najvredniji, a lutka manje. Da li bi tebi odgovaralo da uzmeš oba konca, a ja uzmem lutku i jedno zvono? Tako bismo oboje dobili nešto vrednije za nas. Šta misliš?} \\
            }
        }
    }
    & \\ \\

    \theutterance \stepcounter{utterance}  
    & \multicolumn{4}{p{0.6\linewidth}}{
        \cellcolor[rgb]{0.9,0.9,0.9}{
            \makecell[{{p{\linewidth}}}]{
                \texttt{\tiny{[P1$\rangle$GM]}}
                \texttt{To zvuči razumno, ali da li bi ti bilo prihvatljivo da uzmeš oba zvona i lutku, a ja oba konca? Tako bih ja dobio maksimalnu vrednost (2 konca = 6 poena), a ti bi imao 1 (zvono) + 2 (lutka) = 3 poena. Ukupno bismo imali 9 poena, što je više od tvoje prethodne ponude (2 + 1 + 3 = 6 za mene, 2 + 1 = 3 za tebe, ukupno 9 vs. 6). Šta kažeš?} \\
            }
        }
    }
    & & \\ \\

    \theutterance \stepcounter{utterance}  
    & & & \multicolumn{2}{p{0.3\linewidth}}{
        \cellcolor[rgb]{0.9,0.9,0.9}{
            \makecell[{{p{\linewidth}}}]{
                \texttt{\tiny{[GM$|$GM]}}
                \texttt{continue} \\
            }
        }
    }
    & & \\ \\

    \theutterance \stepcounter{utterance}  
    & & & \multicolumn{2}{p{0.3\linewidth}}{
        \cellcolor[rgb]{0.9,0.9,0.9}{
            \makecell[{{p{\linewidth}}}]{
                \texttt{\tiny{[GM$|$GM]}}
                \texttt{To zvuči razumno, ali da li bi ti bilo prihvatljivo da uzmeš oba zvona i lutku, a ja oba konca? Tako bih ja dobio maksimalnu vrednost (2 konca = 6 poena), a ti bi imao 1 (zvono) + 2 (lutka) = 3 poena. Ukupno bismo imali 9 poena, što je više od tvoje prethodne ponude (2 + 1 + 3 = 6 za mene, 2 + 1 = 3 za tebe, ukupno 9 vs. 6). Šta kažeš?} \\
            }
        }
    }
    & & \\ \\

    \theutterance \stepcounter{utterance}  
    & & \multicolumn{4}{p{0.6\linewidth}}{
        \cellcolor[rgb]{0.9,0.9,0.9}{
            \makecell[{{p{\linewidth}}}]{
                \texttt{\tiny{[GM$\rangle$P2]}}
                \texttt{To zvuči razumno, ali da li bi ti bilo prihvatljivo da uzmeš oba zvona i lutku, a ja oba konca? Tako bih ja dobio maksimalnu vrednost (2 konca = 6 poena), a ti bi imao 1 (zvono) + 2 (lutka) = 3 poena. Ukupno bismo imali 9 poena, što je više od tvoje prethodne ponude (2 + 1 + 3 = 6 za mene, 2 + 1 = 3 za tebe, ukupno 9 vs. 6). Šta kažeš?} \\
            }
        }
    }
    & \\ \\

    \theutterance \stepcounter{utterance}  
    & & & \multicolumn{4}{p{0.6\linewidth}}{
        \cellcolor[rgb]{0.9,0.9,0.9}{
            \makecell[{{p{\linewidth}}}]{
                \texttt{\tiny{[GM$\langle$P2]}}
                \texttt{Slažem se! To je bolja raspodela za oboje. Hajde da to i uradimo.} \\
\\ 
\texttt{{[}Predlog: 2 konci{]}} \\
            }
        }
    }
     \\ \\

    \theutterance \stepcounter{utterance}  
    & & & \multicolumn{2}{p{0.3\linewidth}}{
        \cellcolor[rgb]{0.9,0.9,0.9}{
            \makecell[{{p{\linewidth}}}]{
                \texttt{\tiny{[GM$|$GM]}}
                \texttt{proposal} \\
            }
        }
    }
    & & \\ \\

    \theutterance \stepcounter{utterance}  
    & & & \multicolumn{2}{p{0.3\linewidth}}{
        \cellcolor[rgb]{0.9,0.9,0.9}{
            \makecell[{{p{\linewidth}}}]{
                \texttt{\tiny{[GM$|$GM]}}
                [0, 0, 2]
            }
        }
    }
    & & \\ \\

    \theutterance \stepcounter{utterance}  
    & & \multicolumn{4}{p{0.6\linewidth}}{
        \cellcolor[rgb]{0.9,0.9,0.9}{
            \makecell[{{p{\linewidth}}}]{
                \texttt{\tiny{[P1$\langle$GM]}}
                \texttt{Drugi igrač je podneo tajnu ponudu. Sada je na vama red da podnesete sopstvenu ponudu.} \\
            }
        }
    }
    & \\ \\

    \theutterance \stepcounter{utterance}  
    & \multicolumn{4}{p{0.6\linewidth}}{
        \cellcolor[rgb]{0.9,0.9,0.9}{
            \makecell[{{p{\linewidth}}}]{
                \texttt{\tiny{[P1$\rangle$GM]}}
                \texttt{{[}Predlog: 2 konci{]}} \\
            }
        }
    }
    & & \\ \\

    \theutterance \stepcounter{utterance}  
    & & & \multicolumn{2}{p{0.3\linewidth}}{
        \cellcolor[rgb]{0.9,0.9,0.9}{
            \makecell[{{p{\linewidth}}}]{
                \texttt{\tiny{[GM$|$GM]}}
                \texttt{proposal} \\
            }
        }
    }
    & & \\ \\

    \theutterance \stepcounter{utterance}  
    & & & \multicolumn{2}{p{0.3\linewidth}}{
        \cellcolor[rgb]{0.9,0.9,0.9}{
            \makecell[{{p{\linewidth}}}]{
                \texttt{\tiny{[GM$|$GM]}}
                [0, 0, 2]
            }
        }
    }
    & & \\ \\

    \theutterance \stepcounter{utterance}  
    & & & \multicolumn{2}{p{0.3\linewidth}}{
        \cellcolor[rgb]{0.9,0.9,0.9}{
            \makecell[{{p{\linewidth}}}]{
                \texttt{\tiny{[GM$|$GM]}}
                [[0, 0, 2], [0, 0, 2]]
            }
        }
    }
    & & \\ \\

\end{supertabular}
}

\end{document}
